\documentclass[a4paper, twoside, 12pt]{report}

\usepackage{titlesec}
\usepackage{polski}
\usepackage[english, polish]{babel}
\usepackage{graphicx}
\usepackage[hidelinks]{hyperref}
\usepackage{pdfpages}
\graphicspath{ {./images/} }
\usepackage[utf8]{inputenc}
\usepackage[margin=25mm]{geometry}
\usepackage{indentfirst}
\linespread{1}
\titleformat*{\section}{\fontsize{14pt}{2}\bfseries}
\titleformat*{\subsection}{\fontsize{13pt}{2}\bfseries}
\titleformat*{\subsubsection}{\fontsize{13pt}{2}\bfseries}
\usepackage[T1]{fontenc}
% \usepackage{tgtermes}

\newenvironment{abstractpage}
  {\vspace*{\fill}\thispagestyle{empty}}
    {\vfill}
    \renewenvironment{abstract}[1]
      {\bigskip\selectlanguage{#1}%
             \begin{center}\bfseries\abstractname\end{center}}
           {\par\bigskip}

\begin{document}

% \includepdf{stronatytulowa.pdf}


\begin{abstractpage}
\begin{abstract}{polish}
\end{abstract}

\begin{abstract}{english}
\end{abstract}
\end{abstractpage}

\tableofcontents

\listoffigures

\chapter{Cel i~zakres pracy}
    Celem niniejszej pracy jest przedstawienie zaprojektowanego protokołu kanału
    ukrytego, działającego w~sieciach komputerowych i~ opartego o~ukrywanie wiadomości za pomocą długości pakietów UDP
    zawierających wiadomości podstawowe. w~ramach tego celu zostaną wykonane:
    \begin{itemize}
        \item prezentacja steganografii, nauki o~ukrywaniu wiadomości w~różnego rodzaju nośnikach,
            oraz steganoanalizy, nauki o~wykrywaniu, wraz z ewentualnym niszczeniem, kanałów
            ukrytej komunikacji, wraz z najczęściej używanymi pojęciami
        \item przedstawienie przykładowych, innych niż zaproponowany w~pracy,
            sposobów organizacji kanałów ukrytych, z zakresu steganografii sieciowej,
            z wykorzystaniem różnego rodzaju protokołów komunikacyjnych
        \item zaprezentowanie projektu protokołu kanału ukrytego, ukrywającego
            wiadomości ukryte za pomocą długości pakietów protokołu UDP, wraz
            z opisem parametrów otrzymanego protokołu, oraz wyjaśnieniem pojęć
            użytych do jego opisu
        \item przeprowadzenie badań czynników wpływających na średnie opóźnienie
            w~odbiorze wiadomości podstawowych i~ukrytych, wynikających z wykorzystania
            przedstawionego kanału ukrytego, wraz z analizą wniosków
            z nich płynących
        \item przeprowadzenie analizy wykrywalności zaproponowanego kanału ukrytego,
            wraz z zaproponowaniem sposobów jej zmniejszenia
    \end{itemize}

\chapter{Wstęp teoretyczny}
\clearpage
\addcontentsline{toc}{chapter}{Bibliografia}
\bibliographystyle{plain}
%\bibliography{bibliografia}


\end{document}
